\begin{abstract}
Backtracking regular-expression engines expose a gap between two familiar views of
regular expressions.  Automata theory explains nondeterminism through accepting
runs and their degree of ambiguity, while production engines explain matching
through an ordered backtracking execution that may explore many failing and
successful alternatives before returning the highest-priority match.  The gap is
especially visible in JavaScript, whose semantics is specified operationally and
whose performance pathologies are central to Regular Expression Denial of Service
(ReDoS).

This paper develops a mechanized bridge between these views.  We use
position automata as a common representation and relate accepting runs of a
Glushkov-style automaton to leaves in a JavaScript-style backtracking trace tree.
The resulting correspondence explains when the size and growth of backtracking
search are controlled by classical ambiguity notions: unambiguous automata induce
at most one successful trace, finite ambiguity yields a uniform bound on
successful alternatives, and polynomial or exponential ambiguity gives a semantic
classification of families of inputs that may drive ReDoS behavior.  The
development is being mechanized in Coq.  The current artifact formalizes regular
and positioned syntax, Kleene semantics, position-set constructions, executable
automata, and basic ambiguity definitions; the remaining proof obligations are
isolated as the trace/automaton correspondence and the Weber--Seidl ambiguity
classification layer.
\end{abstract}
